%------------------------------------------------------------------------------%
\begin{question}
  Resolva o sistema linear pelo método de Gauss-Jordan
%  \[
%    \systeme[sort={x,y,z,w,v}]{
%       x + 2y -  z +  w +  v =  3,
%      2x + 5y +  z + 2w +  v =  8,
%      -x -  y + 2z -  w + 2v = -1,
%            y +  z      + 3v =  2
%    }
%  \]
  \[
    \left\{\begin{array}{rrrrrr}
       x & + 2y & -  z & +  w & +  v & = \hpm 3 \\
      2x & + 5y & +  z & + 2w & +  v & = \hpm 8 \\
      -x & -  y & + 2z & -  w & + 2v & =     -1 \\
         &    y & +  z &      & + 3v & = \hpm 2
    \end{array}\right.
  \]
\end{question}

%------------------------------------------------------------------------------%
\begin{resolution}{Matriz Aumentada}
  Matriz aumentada
  \[
    \begin{amatrix}{5}
      1  & 2  & -1 & 1  & 1 & 3  \\
      2  & 5  & 1  & 2  & 1 & 8  \\
      -1 & -1 & 2  & -1 & 2 & -1 \\
      0  & 1  & 1  & 0  & 3 & 2
    \end{amatrix}
  \]
\end{resolution}

%------------------------------------------------------------------------------%
\begin{resolution}{Escalonamento do Sistema}
\begin{columns}
\column{0.45\linewidth}
  \[
    \begin{amatrix}{5}
      1  & 2  & -1 & 1  & 1 & 3  \\
      2  & 5  & 1  & 2  & 1 & 8  \\
      -1 & -1 & 2  & -1 & 2 & -1 \\
      0  & 1  & 1  & 0  & 3 & 2
    \end{amatrix}
  \]
  \pause
\column{0.58\linewidth}
  Eliminar os elementos abaixo do primeiro pivô
  \\ \pause
  \(
    L_2 \leftarrow L_2 - 2 L_1
  \)
  \\ \pause
  \(
    L_3 \leftarrow L_3 + L_1
  \)
\end{columns}
\vspace{-\baselineskip}
\[
  \pause
  \hspace{-2.5em}
  L'_2
  \pause
  \;=\;
  \begin{amatrix}{5}
    2  & 5  & 1  & 2  & 1 & 8
  \end{amatrix}
  -2
  \begin{amatrix}{5}
    1  & 2  & -1 & 1  & 1 & 3
  \end{amatrix}
  \pause
  \;=\;
  \begin{amatrix}{5}
    0 & 1 & 3  & 0 & -1 & 2
  \end{amatrix}
\]
\[
  \pause
  \hspace{-2.5em}
  L'_3
  \pause
  \;=\;
  \begin{amatrix}{5}
    -1 & -1 & 2  & -1 & 2 & -1
  \end{amatrix}
  +
  \begin{amatrix}{5}
    1  & 2  & -1 & 1  & 1 & 3
  \end{amatrix}
  \pause
  \;=\;
  \begin{amatrix}{5}
    0 & 1 & 1  & 0 & 3  & 2
  \end{amatrix}
\]
\pause
\[
  \begin{amatrix}{5}
    1 & 2 & -1 & 1 & 1  & 3 \\
    0 & 1 & 3  & 0 & -1 & 2 \\
    0 & 1 & 1  & 0 & 3  & 2 \\
    0 & 1 & 1  & 0 & 3  & 2
  \end{amatrix}
\]
\end{resolution}

%------------------------------------------------------------------------------%
\begin{resolution}{Escalonamento do Sistema}
\begin{columns}
\column{0.45\linewidth}
  \[
    \begin{amatrix}{5}
      1 & 2 & -1 & 1 & 1  & 3 \\
      0 & 1 & 3  & 0 & -1 & 2 \\
      0 & 1 & 1  & 0 & 3  & 2 \\
      0 & 1 & 1  & 0 & 3  & 2
    \end{amatrix}
  \]
  \pause
\column{0.58\linewidth}
  Eliminar os elementos abaixo do segundo pivô
  \\ \pause
  \(
    L_3 \leftarrow L_3 - L_2
  \)
  \\ \pause
  \(
    L_4 \leftarrow L_4 - L_2
  \)
\end{columns}
\vspace{-\baselineskip}
\[
  \pause
  \hspace{-1.5em}
  L'_3
  \pause
  \;=\;
  \begin{amatrix}{5}
    0 & 1 & 1  & 0 & 3  & 2
  \end{amatrix}
  -
  \begin{amatrix}{5}
    0 & 1 & 3  & 0 & -1 & 2
  \end{amatrix}
  \pause
  \;=\;
  \begin{amatrix}{5}
    0 & 0 & -2 & 0 & 4  & 0
  \end{amatrix}
\]
\[
  \pause
  \hspace{-1.5em}
  L'_4
  \pause
  \;=\;
  \begin{amatrix}{5}
    0 & 1 & 1  & 0 & 3  & 2
  \end{amatrix}
  -
  \begin{amatrix}{5}
    0 & 1 & 3  & 0 & -1 & 2
  \end{amatrix}
  \pause
  \;=\;
  \begin{amatrix}{5}
    0 & 0 & -2 & 0 & 4  & 0
  \end{amatrix}
\]
\pause
\[
  \begin{amatrix}{5}
    1 & 2 & -1 & 1 & 1  & 3 \\
    0 & 1 & 3  & 0 & -1 & 2 \\
    0 & 0 & -2 & 0 & 4  & 0 \\
    0 & 0 & -2 & 0 & 4  & 0
  \end{amatrix}
\]
\end{resolution}

%------------------------------------------------------------------------------%
\begin{resolution}{Escalonamento do Sistema}
\begin{columns}
\column{0.45\linewidth}
  \[
    \begin{amatrix}{5}
      1 & 2 & -1 & 1 & 1  & 3 \\
      0 & 1 & 3  & 0 & -1 & 2 \\
      0 & 0 & -2 & 0 & 4  & 0 \\
      0 & 0 & -2 & 0 & 4  & 0
    \end{amatrix}
  \]
  \pause
\column{0.58\linewidth}
  Tornar o terceiro pivô igual a 1
  \\ \pause
  \(
    L_3 \leftarrow -\frac{1}{2}L_3
  \)
\end{columns}
\vspace{-\baselineskip}
\[
  \pause
  \hspace{-1.5em}
  L'_3
  \pause
  \;=\;
  -\frac{1}{2}
  \begin{amatrix}{5}
    0 & 0 & -2 & 0 & 4  & 0
  \end{amatrix}
  \pause
  \;=\;
  \begin{amatrix}{5}
    0 & 0 & 1  & 0 & -2 & 0
  \end{amatrix}
\]
\pause
\[
  \begin{amatrix}{5}
    1 & 2 & -1 & 1 & 1  & 3 \\
    0 & 1 & 3  & 0 & -1 & 2 \\
    0 & 0 & 1  & 0 & -2 & 0 \\
    0 & 0 & -2 & 0 & 4  & 0
  \end{amatrix}
\]
\end{resolution}

%------------------------------------------------------------------------------%
\begin{resolution}{Escalonamento do Sistema}
\begin{columns}
\column{0.4\linewidth}
  \[
    \begin{amatrix}{5}
      1 & 2 & -1 & 1 & 1  & 3 \\
      0 & 1 & 3  & 0 & -1 & 2 \\
      0 & 0 & 1  & 0 & -2 & 0 \\
      0 & 0 & -2 & 0 & 4  & 0
    \end{amatrix}
  \]
  \pause
\column{0.6\linewidth}
  Eliminar os elementos abaixo do terceiro pivô
  \\ \pause
  \(
    L_4 \leftarrow L_4 +2 L_3
  \)
\end{columns}
\vspace{-\baselineskip}
\[
  \pause
  \hspace{-1.5em}
  L'_4
  \pause
  \;=\;
  \begin{amatrix}{5}
    0 & 0 & -2 & 0 & 4  & 0
  \end{amatrix}
  +2
  \begin{amatrix}{5}
    0 & 0 & 1  & 0 & -2 & 0
  \end{amatrix}
  \pause
  \;=\;
  \begin{amatrix}{5}
    0 & 0 & 0 & 0 & 0 & 0
  \end{amatrix}
\]
\pause
\[
  \begin{amatrix}{5}
    1 & 2 & -1 & 1 & 1  & 3 \\
    0 & 1 & 3  & 0 & -1 & 2 \\
    0 & 0 & 1  & 0 & -2 & 0 \\
    0 & 0 & 0  & 0 &  0 & 0
  \end{amatrix}
\]
\end{resolution}

%------------------------------------------------------------------------------%
\begin{resolution}{Escalonamento do Sistema}
\begin{columns}
\column{0.4\linewidth}
  \[
    \begin{amatrix}{5}
      1 & 2 & -1 & 1 & 1  & 3 \\
      0 & 1 & 3  & 0 & -1 & 2 \\
      0 & 0 & 1  & 0 & -2 & 0 \\
      0 & 0 & 0  & 0 &  0 & 0
    \end{amatrix}
  \]
  \pause
\column{0.6\linewidth}
  Eliminar os elementos acima do terceiro pivô
  \\ \pause
  \(
    L_2 \leftarrow L_2 -3 L_3
  \)
  \\ \pause
  \(
    L_1 \leftarrow L_1 + L_3
  \)
\end{columns}
\vspace{-\baselineskip}
\[
  \pause
  \hspace{-1.5em}
  L'_2
  \pause
  \;=\;
  \begin{amatrix}{5}
    0 & 1 & 3  & 0 & -1 & 2
  \end{amatrix}
  -3
  \begin{amatrix}{5}
    0 & 0 & 1  & 0 & -2 & 0
  \end{amatrix}
  \pause
  \;=\;
  \begin{amatrix}{5}
    0 & 1 & 0 & 0 & 5  & 2
  \end{amatrix}
\]
\[
  \pause
  \hspace{-1.5em}
  L'_1
  \pause
  \;=\;
  \begin{amatrix}{5}
    1 & 2 & -1 & 1 & 1  & 3
  \end{amatrix}
  +
  \begin{amatrix}{5}
    0 & 0 & 1  & 0 & -2 & 0
  \end{amatrix}
  \pause
  \;=\;
  \begin{amatrix}{5}
    1 & 2 & 0 & 1 & -1 & 3
  \end{amatrix}
\]
\pause
\[
  \begin{amatrix}{5}
    1 & 2 & 0 & 1 & -1 & 3 \\
    0 & 1 & 0 & 0 & 5  & 2 \\
    0 & 0 & 1 & 0 & -2 & 0 \\
    0 & 0 & 0 & 0 & 0  & 0
  \end{amatrix}
\]
\end{resolution}

%------------------------------------------------------------------------------%
\begin{resolution}{Escalonamento do Sistema}
\begin{columns}
\column{0.4\linewidth}
  \[
    \begin{amatrix}{5}
      1 & 2 & 0 & 1 & -1 & 3 \\
      0 & 1 & 0 & 0 & 5  & 2 \\
      0 & 0 & 1 & 0 & -2 & 0 \\
      0 & 0 & 0 & 0 & 0  & 0
    \end{amatrix}
  \]
  \pause
\column{0.6\linewidth}
  Eliminar os elementos acima do segundo pivô
  \\ \pause
  \(
    L_1 \leftarrow L_1 - 2 L_2
  \)
\end{columns}
\vspace{-\baselineskip}
\[
  \pause
  \hspace{-1.5em}
  L'_1
  \pause
  \;=\;
  \begin{amatrix}{5}
    1 & 2 & 0 & 1 & -1 & 3
  \end{amatrix}
  -2
  \begin{amatrix}{5}
    0 & 1 & 0 & 0 & 5  & 2
  \end{amatrix}
  \pause
  \;=\;
  \begin{amatrix}{5}
    1 & 0 & 0 & 1 & -11 & -1
  \end{amatrix}
\]
\pause
\[
  \begin{amatrix}{5}
    1 & 0 & 0 & 1 & -11 & -1 \\
    0 & 1 & 0 & 0 & 5   & 2  \\
    0 & 0 & 1 & 0 & -2  & 0  \\
    0 & 0 & 0 & 0 & 0   & 0
  \end{amatrix}
\]
\end{resolution}

%------------------------------------------------------------------------------%
\begin{resolution}{Sistema Linear Reduzido}
%  \[
%    \systeme[sort={x,y,z,w,v}]{
%      x +  w -11v = -1,
%      y + 5v      = 2,
%      z - 2v      = 0
%    }
%  \]
  \[
    \left\{\begin{array}{llllll}
      x &   &   & + w & - 11v & = -1 \\
        & y &   &     & +  5v & = \hpm 2 \\
        &   & z &     & -  2v & = \hpm 0
    \end{array}\right.
  \]
  \pause
  As incógnitas $w$ e $v$ não correspondem a pivôs

  \pause
  portanto são variáveis livres
  \begin{align*}
    w & = s \\
    v & = t
  \end{align*}
  para quaisquer valore reais $t$ e $s$
\end{resolution}

%------------------------------------------------------------------------------%
\begin{resolution}{Solução}
  \[
    \left\{\begin{array}{llllll}
      x &   &   & + w & - 11v & = -1 \\
        & y &   &     & +  5v & = \hpm 2 \\
        &   & z &     & -  2v & = \hpm 0
    \end{array}\right.
  \]
%  \[
%    \systeme[sort={x,y,z,w,v}]{
%      x +  w -11v = -1,
%      y + 5v      = 2,
%      z - 2v      = 0
%    }
%  \]
  \pause
  \vspace{-\baselineskip}
  \begin{columns}[t]
  \column{0.35\linewidth}
    \begin{align*}
      x & = 11v - w -1 \\[-1.5ex]
      y & = 2 - 5v     \\[-1.5ex]
      z & = 2v
    \end{align*}
    \pause
  \column{0.6\linewidth}
    \begin{align*}
      x & = 11t - s - 1 \\[-1.5ex]
      y & = 2 - 5t      \\[-1.5ex]
      z & = 2t          \qquad\qquad\qquad \forall (s, t) \in \R^2 \\[-1.5ex]
      w & = s           \\[-1.5ex]
      v & = t
    \end{align*}
  \end{columns}
\end{resolution}

%------------------------------------------------------------------------------%
\endinput
